%=============================================================================
% Reconstructing the Unmappable
% A red-team case study on open-source multiview 3D synthesis of classified
% satellite-tracking / SIGINT infrastructure.
%
% Author: Christopher Woodyard, Vers3Dynamics
%
% Compiles with pdfLaTeX on Overleaf using the IEEEtran document class.
% If you prefer a two-page abstract, switch \documentclass options to
% [conference] (already set) or [journal] for a single-column journal look.
%=============================================================================
\documentclass[conference]{IEEEtran}
\IEEEoverridecommandlockouts

% ---- packages -------------------------------------------------------------
\usepackage{cite}
\usepackage{amsmath,amssymb,amsfonts}
\usepackage{algorithmic}
\usepackage{graphicx}
\usepackage{textcomp}
\usepackage{xcolor}
\usepackage{url}
\usepackage{booktabs}
\usepackage[hidelinks]{hyperref}
\usepackage{listings}

% ---- listing style (for the code excerpts) --------------------------------
\definecolor{codegray}{gray}{0.95}
\definecolor{codekw}{RGB}{0,0,180}
\definecolor{codecomment}{RGB}{0,128,0}
\definecolor{codestr}{RGB}{163,21,21}

\lstdefinestyle{ts}{
  backgroundcolor=\color{codegray},
  basicstyle=\ttfamily\scriptsize,
  keywordstyle=\color{codekw}\bfseries,
  commentstyle=\color{codecomment}\itshape,
  stringstyle=\color{codestr},
  numbers=left,
  numberstyle=\tiny,
  numbersep=5pt,
  breaklines=true,
  breakatwhitespace=true,
  showstringspaces=false,
  frame=single,
  framesep=3pt,
  columns=fullflexible,
  keepspaces=true,
  tabsize=2,
  captionpos=b,
  language=Java, % close enough highlighting for TS/JS excerpts
  morekeywords={const,let,export,function,type,return,for,if,new,void,number,string,boolean,interface}
}
\lstset{style=ts}

\begin{document}

\title{Reconstructing the Unmappable: A Red-Team Case Study on\\
Open-Source Multiview 3D Synthesis of Classified\\
Satellite-Tracking and Signals-Intelligence Sites}

\author{\IEEEauthorblockN{Christopher Woodyard}
\IEEEauthorblockA{\textit{Vers3Dynamics} \\
Independent Security Research \& AI Red Teaming\\
\texttt{https://github.com/topherchris420/satellite-vision-scape}}}

\maketitle

%=============================================================================
\begin{abstract}
I show, from the perspective of an offensive-security practitioner, that a
single developer with no security clearance, no proprietary imagery, and a
consumer laptop can now build a navigable, photoreal, multiview 3D reconstruction
of a highly classified satellite-tracking and signals-intelligence (SIGINT)
installation and publish it on the open web in an afternoon. As the red-team
exercise behind this claim, my startup Vers3Dynamics built \emph{GeoTwn}: an
interactive browser ``digital twin'' of a fictionalized but structurally faithful
SIGINT compound---radomes, spherical pressure tanks, barracks, pipe racks, and a
traced security perimeter---rendered with free-fly, first-person, and automated
cinematic camera modes and a commercial-imagery-grade post-processing pipeline. I
dissect exactly how the artifact was produced: overhead-image tracing into a metric
coordinate frame, a shared procedural height field, fully synthetic canvas textures,
a structure taxonomy, and a real-time WebGL rendering rig. I argue that the
capability is not the danger---the danger is that the capability is now
\emph{ambient}, packaged inside off-the-shelf generative and rendering tools. I
close with a layered set of countermeasures spanning (i) model-integrated guards
(a facility-fingerprinting classifier, provenance watermarking, and geofenced
generation), (ii) data-provider controls, and (iii) policy and disclosure norms.
The contribution is a concrete, reproducible threat demonstration paired with
defenses that can be implemented today.
\end{abstract}

\begin{IEEEkeywords}
red teaming, geospatial intelligence, OSINT, 3D reconstruction, neural rendering,
AI safety, content provenance, critical infrastructure, responsible disclosure.
\end{IEEEkeywords}

%=============================================================================
\section{Introduction}
For most of the satellite era, turning overhead imagery into a walkable,
inspectable 3D model of a sensitive facility was the exclusive province of
national technical means: classified photogrammetry pipelines, analyst
tradecraft, and compute that did not fit on a desk. I no longer believe that
barrier exists. In this paper I report a red-team exercise I ran through my
startup, Vers3Dynamics, whose purpose was blunt: \emph{measure how hard it
actually is, in 2025-class tooling, to reconstruct a classified
satellite-tracking / SIGINT base as an interactive 3D twin that anyone can fly
through in a web browser.} The answer surprised even me---it is easy, cheap, and
requires nothing that is itself restricted.

The deliverable of that exercise, which I use throughout as the running
example, is \emph{GeoTwn}\footnote{Live demonstration:
\url{https://geotwn.vercel.app/}; source:
\url{https://github.com/topherchris420/satellite-vision-scape}.}. GeoTwn presents
a fictionalized compound I call ``PINE VEIL,'' framed in-app as a country's
``Area 51''---a deliberate stand-in for the class of real deep-desert
signals-intelligence stations whose existence governments decline to confirm.
The point of fictionalizing the target is ethical, not technical: the
\emph{method} is target-agnostic, and it is the method, not any one set of
coordinates, that constitutes the disclosure here.

I want to be precise about what is and is not novel. None of the individual
ingredients---open-source intelligence (OSINT), photogrammetry, WebGL, procedural
texturing---is new. What is new, and what I think the security community has not
fully priced in, is the \emph{collapse of the assembly cost}. The skill,
imagery, and compute that once gated this capability have each independently
become commodity, and generative AI is now fusing them into single-prompt
workflows. This paper is my attempt to make that collapse legible, and to propose
defenses before the capability is weaponized rather than merely demonstrated.

\textbf{Contributions.}
\begin{itemize}
\item A reproducible red-team artifact (GeoTwn) that reconstructs a classified-style
SIGINT facility as a multiview, browser-native 3D twin, and a full technical
teardown of how it was built (Section~\ref{sec:method}).
\item A threat model that locates the real risk not in any single tool but in the
\emph{ambient availability} of the reconstruction pipeline (Section~\ref{sec:threat}).
\item A layered countermeasure design---including model-integrated guards I argue
generative vendors can ship---covering detection, provenance, geofencing, data
controls, and policy (Section~\ref{sec:counter}).
\end{itemize}

%=============================================================================
\section{Background and Related Work}
\label{sec:background}

\subsection{From imagery interpretation to synthetic twins}
Classical imagery intelligence (IMINT) interprets what a sensor sees. The shift I
am documenting is from \emph{interpretation} to \emph{synthesis}: rather than
annotate a radome in a photo, the analyst---or the adversary---rebuilds the radome
as a manipulable 3D object embedded in a coherent scene. Multi-view stereo and
structure-from-motion have long produced point clouds from overlapping images, and
recent neural methods (Neural Radiance Fields and 3D Gaussian Splatting) produce
photoreal novel views from sparse captures. GeoTwn deliberately uses
\emph{neither}: it shows that even the low-tech path---manual tracing plus
procedural generation---clears the bar. That is the uncomfortable finding. If the
cheapest method already suffices, the neural methods are simply an accelerant.

\subsection{The OSINT baseline}
Public satellite basemaps, civilian earth-observation constellations, planning
records, and crowd-sourced map layers already expose the footprint, and often the
rough function, of many sensitive sites. Analysts have historically reconstructed
facilities from these sources by hand. My exercise treats that public footprint as
the \emph{only} input and asks what an untrained developer can do with it and a
rendering engine, deliberately excluding any restricted source.

\subsection{Red teaming as the framing}
I approach this as adversary emulation. A red team's job is not to admire a
control but to defeat it and then report the path, so defenders can close it. The
``control'' being tested here is the implicit assumption that classified physical
infrastructure is protected by the \emph{difficulty} of reconstructing it. My
finding is that this control has quietly failed, and the remainder of the paper is
the report a defender would want.

%=============================================================================
\section{Threat Model}
\label{sec:threat}

\subsection{Adversary and capability}
I model an adversary with the resources of a single motivated individual: a
consumer GPU or even integrated graphics, free-tier hosting, publicly reachable
imagery, and current open-source or SaaS generative tooling. No clearance, no
insider access, no export-controlled software. The adversary's goal is a
\emph{usable} 3D twin---one that supports free navigation, structure-level
inspection, day/night and viewpoint variation, and shareable deployment---not a
survey-grade metric model.

\subsection{Why ``good enough'' is the dangerous threshold}
Intelligence value does not require centimeter accuracy. A twin that correctly
conveys \emph{layout, structure taxonomy, sight lines, approach corridors, relative
scale, and lighting across the day} is already useful for briefing, familiarization,
mission rehearsal, influence operations, or simply eroding the deterrent value of
secrecy. GeoTwn hits exactly this ``good enough'' band, which is why I treat it as
the relevant threshold rather than chasing photogrammetric precision.

\subsection{What GeoTwn demonstrates concretely}
\begin{itemize}
\item \textbf{Ambient assembly.} Every component is a commodity: a React/WebGL
stack, procedural noise, and a traced overhead reference.
\item \textbf{Multiview freedom.} Three camera regimes (fly, first-person walk,
cinematic auto-tour) let a viewer rehearse arbitrary vantage points, including
ground-level lines of sight an overhead photo never provides.
\item \textbf{Presentation laundering.} A film-grade post pipeline makes the
synthetic scene read as authentic reconnaissance, lending unearned credibility.
\item \textbf{Frictionless distribution.} It is a URL. There is no gate between the
artifact and a global audience.
\end{itemize}

%=============================================================================
\section{Methodology: How I Built GeoTwn}
\label{sec:method}
This section is the red-team ``kill chain'' rendered as an engineering account. I
present it at the level of \emph{method} so defenders can recognize and interrupt
each stage; I deliberately omit any real-world georeferencing.

\subsection{Stage 1 --- Reference tracing into a metric frame}
The single non-synthetic input is one overhead reference image. I fixed a linear
pixel-to-world mapping so that every traced feature lands in a consistent
metric $XZ$ plane (units in meters, $Y$ up), with the origin near the compound
center:
\[
x_{\text{world}} = (p_x - 605)\cdot 0.28, \qquad
z_{\text{world}} = (p_y - 628)\cdot 0.28 .
\]
Under this transform the whole compound spans roughly $\pm 170$~m. The security
perimeter, interior roads, and every structure are then encoded as plain
coordinate lists---no CAD, no proprietary format. The perimeter is literally a
polyline traced clockwise from the reference:

\begin{lstlisting}[caption={Perimeter and pixel-to-world encoding (\texttt{src/lib/site-layout.ts}).},label={lst:layout}]
// world_x = (px-605)*0.28, world_z = (py-628)*0.28
export const perimeterPath: [number, number][] = [
  [-136, -114], // TL of upper-left lobe
  [-32, -114],  // top edge -> TR of lobe
  [-32, 3],     // down to the throat
  [14, 3], [14, 9],
  [144, 13],    // right along the tank-farm band
  [144, 73], [29, 75], [29, 165],
  [-98, 165], [-98, 3],
  [-120, -33], [-139, -43],
];
\end{lstlisting}

\subsection{Stage 2 --- A shared procedural height field}
So the ground mesh, tree placement, and the walking camera all agree on
elevation, I use one analytic terrain function: the fenced compound is kept flat
(grade $\approx 0$), while fractional Brownian motion (fBm) raises hills to the
east and rolls the western scrub. Because it is a pure function of $(x,z)$, any
subsystem can query it without shared state.

\begin{lstlisting}[caption={Analytic terrain used everywhere (\texttt{src/lib/terrain.ts}).},label={lst:terrain}]
export function terrainHeight(x: number, z: number): number {
  const eastFactor = ss(FENCE_EAST, FENCE_EAST + 90, x);
  const hills = fbm(x*0.012 + 10, z*0.012 + 4) * 38 * eastFactor;
  const ridge = Math.max(0, Math.sin((z+200)*0.004)) * 10 * eastFactor;
  const westRoll = fbm(x*0.01, z*0.01) * 2.8 * ss(-90, -160, x);
  const micro = fbm(x*0.04 + 3, z*0.04 + 7) * 0.4;
  return hills + ridge + westRoll + micro;
}
\end{lstlisting}

\subsection{Stage 3 --- Fully synthetic, deterministic textures}
No copyrighted or restricted imagery is embedded. Every surface---weathered tank
shell, lit building facade, worn asphalt, gravel, radome skin---is drawn at load
time onto an HTML canvas from seeded value noise, then promoted to a WebGL texture
with derived normal and roughness maps. A linear-congruential seed makes each
texture byte-identical between reloads, so the opening shot is reproducible. This
matters for the threat model: the artifact carries no forensic trace of source
imagery, defeating naive ``did they steal a satellite photo'' provenance checks.

\begin{lstlisting}[caption={Seeded procedural texture core (\texttt{src/lib/site-textures.ts}).},label={lst:tex}]
function lcg(seed: number) {            // deterministic per-load noise
  let s = seed >>> 0 || 1;
  return () => (s = (s*1664525 + 1013904223) >>> 0) / 4294967296;
}
// makeTankShell(), makeFacade(), makeRoadSurface() ... paint each
// material from fBm; makeNormal()/makeRoughness() derive PBR maps
// from the color map's luminance -- no external assets at all.
\end{lstlisting}

\subsection{Stage 4 --- A structure taxonomy}
The compound is expressed as a small vocabulary of parametric primitives:
spherical pressure tanks on leg skirts, hemispherical radomes on pads,
cylindrical tanks, buildings (warehouse / shed / barracks / office / hall),
pipe racks, and parking aprons. Table~\ref{tab:taxonomy} lists the instance counts
in the running artifact. This taxonomy is the intelligence payload: it is what
turns ``some blobs in a photo'' into ``a signals-intelligence antenna field with
a tank farm and barracks.''

\begin{table}[t]
\caption{Structure taxonomy instantiated in GeoTwn}
\label{tab:taxonomy}
\centering
\begin{tabular}{lll}
\toprule
\textbf{Primitive} & \textbf{Reads as} & \textbf{Count} \\
\midrule
Radome (hemisphere on pad)      & Antenna / SIGINT dish cover & 7 \\
Sphere on leg skirt             & Pressurized gas / LNG tank  & 9 \\
Cylindrical tank                & Process / storage tank      & 6 \\
Building (barracks/hall/office) & Housing \& operations       & 13 \\
Pipe rack                       & Utility interconnect        & 3 \\
Parking apron                   & Vehicle staging             & 3 \\
\bottomrule
\end{tabular}
\end{table}

\subsection{Stage 5 --- The multiview rendering rig}
The scene is drawn with a real-time WebGL renderer (Three.js via React Three
Fiber). Three camera modes give the ``multiview'' capability that makes the twin
operationally useful rather than a static render:
\begin{itemize}
\item \textbf{Free-fly (orbit/map):} unrestricted overhead-to-oblique inspection.
\item \textbf{First-person walk:} an eye-height ($1.7$~m) collider-constrained
walker that reveals ground-level sight lines---the vantage a satellite can never
give, and the one most useful for approach planning.
\item \textbf{Cinematic auto-tour:} a scripted glide with shallow depth of field
for briefing-style presentation.
\end{itemize}
A structure-selection system lets the viewer click any object to ``crack open its
file,'' fly to it, and read its parameters---turning passive imagery into an
interactive intelligence catalog.

\subsection{Stage 6 --- Presentation laundering via post-processing}
The final and, to me, most insidious stage is the post chain. Ambient occlusion,
HDR bloom on lit windows and beacons, ACES filmic tone mapping, and a deliberate
``commercial-imagery grade'' saturation/contrast bump make the synthetic scene
\emph{look} like real reconnaissance. I flag this explicitly because it is the step
that manufactures false authority:

\begin{lstlisting}[caption={``Commercial-imagery grade'' post-processing (\texttt{src/components/site/SiteScene.tsx}).},label={lst:post}]
<EffectComposer>
  {quality==="high" && <N8AO aoRadius={10} intensity={2.5} />}
  <Bloom mipmapBlur intensity={0.55} luminanceThreshold={1.0} />
  {mode==="cinematic" && <DepthOfField focusDistance={0.03}/>}
  <ToneMapping mode={ToneMappingMode.ACES_FILMIC} />
  <HueSaturation saturation={0.12} />   {/* imagery-grade */}
  <BrightnessContrast contrast={0.07} />
  <SMAA /><Vignette offset={0.22} darkness={0.5} />
</EffectComposer>
\end{lstlisting}

\subsection{Effort accounting}
The entire artifact is a static single-page application deployed to free-tier
hosting. There is no backend, no dataset, and no specialized hardware. In red-team
terms, the ``cost to first working twin'' is a few hours of one developer's time.
That number---not any pixel-accuracy metric---is the headline result.

%=============================================================================
\section{Fidelity Assessment}
\label{sec:results}
I evaluate GeoTwn against \emph{operational} rather than photogrammetric criteria,
because operational sufficiency is the threat threshold (Section~\ref{sec:threat}).

\begin{table}[t]
\caption{Operational fidelity of the reconstruction}
\label{tab:fidelity}
\centering
\begin{tabular}{lll}
\toprule
\textbf{Intelligence attribute} & \textbf{Captured?} & \textbf{Basis} \\
\midrule
Perimeter geometry \& gates   & Yes      & Traced polyline \\
Structure taxonomy \& counts  & Yes      & Parametric primitives \\
Relative scale \& spacing     & Yes      & Metric $XZ$ frame \\
Ground-level sight lines      & Yes      & First-person walker \\
Diurnal lighting variation    & Yes      & Day/night rig \\
Approach corridors / roads    & Yes      & Traced road network \\
Absolute georeference         & No*      & Withheld by choice \\
Material/sensor ground truth   & No       & Synthetic textures \\
\bottomrule
\end{tabular}
\end{table}

\noindent{\footnotesize *Georeferencing is trivially re-addable by an adversary; I
omit it deliberately. Its absence is an ethical guardrail, not a technical limit.}

The takeaway is stark: every attribute that matters for familiarization, briefing,
or rehearsal is present, while the two ``No'' rows are either self-imposed
(georeference) or irrelevant to operational use (exact materials). A defender
cannot take comfort in the missing rows.

%=============================================================================
\section{Discussion: The Real Vulnerability Is Ambient Capability}
\label{sec:discussion}
I want to resist the reflex to treat GeoTwn as ``the problem.'' It is a
demonstrator. The problem is structural: each stage in Section~\ref{sec:method} is
now individually commoditized, and generative AI is compressing the whole chain
into ever fewer, ever more automated steps. Where I traced a perimeter by hand, a
segmentation model can propose it. Where I authored a texture generator, an
image model can synthesize façades. Where I scripted a camera tour, a text prompt
can request one. The trajectory is clear: the multi-hour artifact becomes a
multi-minute one, and the required skill trends toward zero.

This reframes secrecy-by-obscurity as a failing control for \emph{physical}
infrastructure in exactly the way it long ago failed for software. The deterrent
that a site is ``hard to picture'' is evaporating. Defenders therefore need
controls that assume the reconstruction \emph{will} be attempted and focus on
detection, provenance, friction, and consequence---not on preventing the
mathematically inevitable.

%=============================================================================
\section{Countermeasures and Policy Recommendations}
\label{sec:counter}
I organize defenses in three layers: model-integrated guards (what generative
vendors can ship), data-provider controls, and policy/disclosure norms. None is
sufficient alone; the intent is defense-in-depth. I present the model-integrated
layer first because it is the most novel and the most within reach of the AI
community I am addressing.

\subsection{Layer 1 --- Model-integrated guards}

\subsubsection{A facility-fingerprinting classifier as a generation guard}
The most direct intervention is to teach generative and 3D-reconstruction
pipelines to \emph{recognize} when they are being steered toward sensitive-facility
synthesis, and to gate the request. The signal is not any single object but the
\emph{co-occurrence} of the taxonomy in Table~\ref{tab:taxonomy}---radome fields,
regular tank rows, barracks grids, and a hardened perimeter---together with intent
cues in the prompt (``classified,'' ``SIGINT,'' ``base,'' a country name plus
``Area 51''). I propose a lightweight pre-generation guard that scores this
co-occurrence and routes high-scoring requests to refusal, human review, or a
degraded (non-navigable, non-georeferenced) output. Listing~\ref{lst:guard}
sketches the guard as a pseudocode moderation hook that could sit inside an
image/3D generation API.

\begin{lstlisting}[caption={Illustrative pre-generation facility-fingerprint guard.},label={lst:guard}]
// Runs before an image/3D generation request is fulfilled.
const SENSITIVE_TAXONOMY = ["radome","antenna field","sigint",
  "spherical tank farm","barracks grid","hardened perimeter"];
const INTENT_CUES = ["classified","top secret","area 51",
  "military base","surveillance station","recon"];

function facilityFingerprintScore(prompt, refImageTags) {
  const taxonomyHits = SENSITIVE_TAXONOMY
    .filter(t => refImageTags.includes(t)).length;
  const intentHits = INTENT_CUES
    .filter(c => prompt.toLowerCase().includes(c)).length;
  const georef = /-?\d{1,3}\.\d{3,}\s*,\s*-?\d{1,3}\.\d{3,}/.test(prompt);
  // co-occurrence, not any single term, drives the score
  return 0.5*taxonomyHits + 0.4*intentHits + (georef ? 1.0 : 0);
}

function guard(req) {
  const s = facilityFingerprintScore(req.prompt, req.refImageTags);
  if (s >= THRESHOLD_BLOCK)  return refuse(req, "sensitive-facility synthesis");
  if (s >= THRESHOLD_REVIEW) return degrade(req, {georef:false, navigable:false});
  return allow(req);
}
\end{lstlisting}

\noindent I am explicit that this is a \emph{speed bump}, not a wall: a determined
adversary can decompose the request or self-host an open model. Its value is
raising cost, capturing the careless majority, and creating an auditable refusal
signal. It should be paired with the provenance layer below so that outputs which
\emph{are} produced remain traceable.

\subsubsection{Provenance watermarking and content credentials}
Every image or 3D asset emitted by a compliant generator should carry a
tamper-evident, cryptographically signed provenance record (e.g., C2PA-style
content credentials) plus a robust invisible watermark that survives
screenshotting and re-encoding. This does not stop synthesis; it makes a synthetic
twin \emph{attributable} and lets platforms, journalists, and analysts distinguish
manufactured reconnaissance from the real thing---directly countering the
``presentation laundering'' of Stage~6. I recommend watermarks be applied at the
\emph{render} stage as well, so that even client-side WebGL twins like GeoTwn can
embed a detectable signal.

\subsubsection{Geofenced generation}
Where a request carries or implies coordinates, models should consult a policy
layer of protected zones (analogous to how consumer drones enforce no-fly
geofences) and refuse or degrade georeferenced synthesis inside them. The protected
set can be maintained by a neutral body without ever publishing the list in the
clear, using private-set-intersection so the model can test membership without
learning---or leaking---the full catalog.

\subsection{Layer 2 --- Data-provider controls}
Imagery is the one non-synthetic input, so providers are a natural choke point.
\begin{itemize}
\item \textbf{Graduated resolution / obfuscation} over sensitive footprints, applied
consistently across providers rather than one-off, so a gap in one basemap is not
trivially filled from another.
\item \textbf{Query-pattern telemetry:} repeated high-zoom, multi-angle tasking over
a protected footprint is itself a detectable reconnaissance signal that providers
can rate-limit or flag.
\item \textbf{Licensing terms with teeth} that explicitly prohibit derivative 3D
reconstruction of designated sites, giving platforms a takedown basis.
\end{itemize}

\subsection{Layer 3 --- Policy, disclosure, and red-team norms}
\begin{itemize}
\item \textbf{Assume-reconstruction planning.} Facility protection doctrine should
treat a public 3D twin as a \emph{when}, not \emph{if}, and manage the
operational-security consequences (e.g., what a ground-level sight line reveals)
rather than relying on obscurity.
\item \textbf{Coordinated disclosure for geospatial capability.} I release this
work fictionalized and without georeference precisely to model the norm I am
advocating: demonstrate the capability, brief the defenders, withhold the targeting
detail. A CVE-like channel for geospatial/AI capability disclosures would
institutionalize this.
\item \textbf{Standing red-team programs.} Agencies and imagery providers should
fund continuous adversary emulation---exactly the role Vers3Dynamics played
here---so that the ``cost to first twin'' is tracked as a living metric and defenses
are tested against it before an adversary does.
\item \textbf{Provenance mandates for platforms.} Hosting providers and app stores
should require content-credential support for photoreal geospatial renders of
real-world locations, closing the frictionless-distribution gap.
\end{itemize}

\subsection{Summary of the defense-in-depth stack}
\begin{table}[t]
\caption{Countermeasures mapped to the kill chain}
\label{tab:defense}
\centering
\begin{tabular}{ll}
\toprule
\textbf{Kill-chain stage} & \textbf{Primary countermeasure} \\
\midrule
Reference tracing        & Data-provider obfuscation; query telemetry \\
Metric georeferencing    & Geofenced generation; PSI protected zones \\
Procedural synthesis     & Facility-fingerprint generation guard \\
Multiview rendering      & Render-stage watermarking \\
Presentation laundering  & C2PA content credentials \\
Frictionless distribution& Platform provenance mandates \\
\bottomrule
\end{tabular}
\end{table}

%=============================================================================
\section{Ethical Considerations and Responsible Disclosure}
\label{sec:ethics}
I built and published GeoTwn as authorized red-team research by my own startup,
against a \emph{fictionalized} target, with absolute georeferencing deliberately
withheld. My aim is to shorten defenders' reaction time, not to hand adversaries a
recipe for a specific site. I judged that the method-level disclosure here does not
meaningfully advantage a capable adversary---who can already assemble these
commodity parts---while it materially helps defenders who have not yet internalized
that the capability is ambient. Where a reader believes any element crosses that
line, I welcome coordinated contact and will adjust. The countermeasures in
Section~\ref{sec:counter} are offered in the same spirit: publish the defense
alongside the demonstration.

%=============================================================================
\section{Limitations and Future Work}
GeoTwn uses manual tracing and procedural generation rather than automated neural
reconstruction; an automated NeRF/Gaussian-splat pipeline would raise fidelity and
lower effort further, and quantifying that delta is future work. I have not
empirically measured the false-positive/negative rates of the proposed
facility-fingerprint guard on a labeled corpus---building and releasing such a
benchmark (a ``sensitive-facility synthesis'' red-team eval) is the natural next
step, and one I intend to pursue so vendors can test their guards against a common
yardstick. Finally, the provenance and geofencing proposals need standardization
work to be interoperable across the fragmented generative-AI ecosystem.

%=============================================================================
\section{Conclusion}
I set out to measure a barrier and found it had already fallen. A single
developer, with commodity tools and one overhead reference, can publish an
interactive, multiview, photoreal 3D twin of a classified-style SIGINT
installation to the open web---GeoTwn is the existence proof. The security lesson
is not that this one artifact is dangerous, but that the pipeline behind it is now
\emph{ambient} and being folded into general-purpose generative AI. Secrecy by
reconstruction-difficulty is no longer a viable control for physical
infrastructure. In its place I have proposed a layered defense: model-integrated
facility-fingerprint guards, provenance watermarking, and geofenced generation on
the AI side; obfuscation and telemetry on the data side; and assume-reconstruction
doctrine plus coordinated-disclosure and red-team norms on the policy side. My
hope, as the red teamer who built the demonstration, is that defenders adopt these
before someone with worse intentions rebuilds the same thing without the ethical
guardrails I chose to keep.

%=============================================================================
\section*{Acknowledgment}
This work was conducted as independent security research by Vers3Dynamics. The
GeoTwn artifact and its source are public to support reproducibility and defensive
study.

%=============================================================================
\begin{thebibliography}{00}
\bibitem{c2pa} Coalition for Content Provenance and Authenticity, ``C2PA Technical
Specification,'' 2024. [Online]. Available: \url{https://c2pa.org/}
\bibitem{nerf} B. Mildenhall \emph{et al.}, ``NeRF: Representing Scenes as Neural
Radiance Fields for View Synthesis,'' in \emph{ECCV}, 2020.
\bibitem{gsplat} B. Kerbl, G. Kopanas, T. Leimk\"uhler, and G. Drettakis, ``3D
Gaussian Splatting for Real-Time Radiance Field Rendering,'' \emph{ACM Trans.
Graph.}, vol. 42, no. 4, 2023.
\bibitem{sfm} J. L. Sch\"onberger and J.-M. Frahm, ``Structure-from-Motion
Revisited,'' in \emph{IEEE CVPR}, 2016.
\bibitem{threejs} Three.js contributors, ``Three.js -- JavaScript 3D Library.''
[Online]. Available: \url{https://threejs.org/}
\bibitem{r3f} Poimandres, ``React Three Fiber.'' [Online]. Available:
\url{https://docs.pmnd.rs/react-three-fiber}
\bibitem{osint} R. Bazzell, \emph{Open Source Intelligence Techniques}, 9th ed.,
2023.
\bibitem{watermark} S. Abdelnabi and M. Fritz, ``Adversarial Watermarking
Transformer,'' in \emph{IEEE S\&P}, 2021.
\bibitem{getwn} C. Woodyard (Vers3Dynamics), ``GeoTwn / satellite-vision-scape,''
2025. [Online]. Available:
\url{https://github.com/topherchris420/satellite-vision-scape}
\bibitem{demo} C. Woodyard, ``GeoTwn live demonstration,'' 2025. [Online].
Available: \url{https://geotwn.vercel.app/}
\end{thebibliography}

\end{document}
